\chapter{Il tempo, misurato dall'anno zero}
\label{cap:tempo_ghiacci_clima}

Per capire dove cercare tracce di frequentazione preistorica sulle Alpi venete, occorre prima sapere quando quelle Alpi sono diventate abitabili, e come sono cambiate nel corso dei millenni. Il clima non è stato uno sfondo immutabile: ha aperto e chiuso territori, spostato il limite delle nevi, fatto avanzare e arretrare le foreste. Questo capitolo ricostruisce quella storia fase per fase, con le date e le fonti che la sostengono, non come premessa generica ma come base di lavoro per tutto ciò che segue.

\section{Una nota sul modo di contare gli anni}

Gli specialisti che studiano questo periodo usano convenzioni di datazione poco intuitive per chi non lavora nel settore: sigle come \termine{BP} (\textit{before present}, ``prima del presente''), dove il presente è fissato per convenzione all'anno 1950, oppure \termine{cal BP}, la stessa scala corretta con le calibrazioni necessarie a compensare le irregolarità del metodo di datazione al radiocarbonio, spesso abbreviata in \termine{ka cal BP} quando espressa in migliaia di anni. Sono strumenti utili per chi deve confrontare datazioni di laboratorio, ma poco maneggevoli per chiunque altro.

In questo libro tutte le date sono espresse in anni prima dell'anno zero del calendario astronomico, quello che si apprende a scuola come ``anno 1 avanti Cristo''. La conversione dalla scala cal BP è semplice: si sottraggono circa 1.950 anni, con un arrotondamento al centinaio o al migliaio più vicino, coerente con i margini di incertezza delle datazioni originali. La tabella \ref{tab:cronologia} raccoglie tutte le fasi discusse in questo capitolo; chi consultasse la letteratura originale, che usa cal BP, può ritrovare il valore originale aggiungendo ~1.950 al numero riportato qui.

\section{Il grande ghiaccio e il suo ritiro}

Circa diciannove mila anni prima dell'anno zero, le Alpi attraversavano la fase più fredda dell'ultima grande glaciazione, quella che i geologi chiamano Ultimo Massimo Glaciale. Le valli principali erano occupate da lingue di ghiaccio che scendevano fin sotto le montagne, riempiendo gli spazi oggi occupati dai grandi laghi prealpini e dalle pianure pedemontane. In queste condizioni, l'idea stessa di una presenza umana stabile in quota non ha senso: il territorio montano era semplicemente inospitale, coperto di ghiaccio o esposto a condizioni periglaciali proibitive.

Il periodo che segue, definito Tardoglaciale, è convenzionalmente fatto iniziare a diciassette mila anni prima dell'anno zero e concluso a circa novemilacinquecento, secondo la sintesi più organica disponibile per l'area alpina italiana, quella di Cesare Ravazzi, Marco Peresani, Roberta Pini ed Elisa Vescovi \citep{ravazzi2007}. Dopo il ritiro dei grandi ghiacciai dagli anfiteatri pedemontani e dai bacini dei grandi laghi prealpini, gran parte delle avanzate glaciali successive si concentra nella prima parte di questo lungo periodo, prima dell'interstadio caldo di cui si parlerà a breve: la vegetazione forestale resta allora confinata alle aree pedemontane e alla fascia esterna delle Prealpi, senza penetrare nelle valli interne, ancora soggette a un'intensa evoluzione paraglaciale, cioè al riassestamento dei versanti liberati dal ghiaccio, con frane, colate e instabilità diffusa.

All'interno di questa prima fase fredda, i geologi riconoscono un'oscillazione più breve, databile tra quindicimila e tredicimilaseicento anni prima dell'anno zero circa, nota come oscillazione di Ragogna \citep{ravazzi2005}, durante la quale la copertura forestale, già timidamente avviata, subisce una battuta d'arresto legata a un episodio di raffreddamento, lo stesso che nell'Europa centro-settentrionale è noto come Heinrich Stadial 1, e che nelle Alpi provoca anche l'avanzata glaciale detta stadio di Gschnitz.

\section{Il primo grande miglioramento climatico}

Un cambiamento più rapido e più significativo comincia in un intervallo compreso tra dodicimilasettecento e dodicimilatrecento anni prima dell'anno zero, quando le condizioni climatiche migliorano sensibilmente per un periodo di alcuni secoli: è l'interstadio noto come Bølling-Allerød. È in questa fase che le foreste avanzano rapidamente in quota, arrivando a colonizzare i versanti fino a milleseicento, millesettecento metri sul livello del mare, e la produttività dei sistemi lacustri e palustri aumenta sensibilmente, tanto che la formazione di torba comincia a manifestarsi fino a quote di millesettecento metri.

Il dato più importante di questa fase, per il tema di questo libro, è che è proprio in questo momento che compaiono le prime tracce sicure di frequentazione umana stabile nella fascia montana delle Alpi orientali: è la colonizzazione epigravettiana delle Prealpi e delle Dolomiti, con un numero di siti che aumenta e si differenzia per quota e per funzione rispetto al periodo precedente.

La seconda metà di questo interstadio caldo è marcata dall'espansione di latifoglie termofile, il cui apice coincide con la deposizione di un sottile strato di cenere vulcanica, la Laacher See Tephra, un evento databile con precisione tra undicimila e diecimilanovecento anni prima dell'anno zero e usato dai geologi come livello guida per sincronizzare le sequenze sedimentarie in tutta l'Europa centrale, perché anticipa di circa due secoli l'inizio del successivo raffreddamento.

\section{Un rovescio temporaneo}

Il miglioramento climatico non fu definitivo. Il Dryas Recente, il raffreddamento che segue l'interstadio caldo, coincide quasi esattamente con l'arrivo della Laacher See Tephra, a conferma che l'inizio di questa fase fredda fu sincrono tra il nord Italia e l'Europa centrale. Durante il Dryas Recente, la pianura padana conobbe un clima continentale secco, prossimo al regime semiarido, con forte contrasto stagionale.

Nella seconda metà di questo periodo freddo comincia una fase complessa di avanzata glaciale, documentata da tre distinti episodi noti come stadi di Egesen, Kartell e Kromer, che si protrae ben oltre la fine convenzionale del Dryas Recente, interessando anche i primi due millenni dell'Olocene. Uno studio sulla Val Cantone di Dosdé, nelle Alpi Retiche centrali, ha permesso di datare direttamente le morene della fase Egesen a novemilaottocento anni prima dell'anno zero circa, con un margine di incertezza di circa millecento anni: un dato che dà concretezza a una fase altrimenti descritta solo in termini relativi.

È interessante notare, e va detto con onestà, che i dati archeologici attualmente disponibili non permettono di stabilire con chiarezza quale effetto abbia avuto questo peggioramento climatico sulle popolazioni che nel frattempo avevano cominciato a frequentare la montagna: gli stessi autori della sintesi qui seguita dichiarano esplicitamente che, allo stato attuale, i dati archeologici non sono in grado di rilevare l'impatto del Dryas Recente sulle dinamiche insediative \citep{ravazzi2007}. Non è chiaro se gli insediamenti di quota siano stati abbandonati, ridotti, o se siano proseguiti pressoché inalterati. È uno di quei punti in cui la ricerca attuale non ha ancora una risposta definitiva, e l'onestà del metodo seguito in questo libro impone di dirlo esplicitamente, invece di colmare il vuoto con una supposizione presentata come fatto accertato.

\section{Un caso locale: il Cansiglio}

Un dato di particolare interesse per il territorio di questo libro viene proprio dalle Prealpi venete, dal bacino di Palughetto sull'altopiano del Cansiglio, dove una successione paleobotanica e archeologica ha permesso di ricostruire nel dettaglio la transizione tra Tardoglaciale e primo Olocene in un contesto prealpino orientale \citep{avigliano2000}. È uno dei pochi punti del territorio veneto per cui esiste una sequenza climatica e vegetazionale ricostruita con continuità stratigrafica diretta, non per estrapolazione da siti più lontani: un promemoria del fatto che, come si vedrà nell'ultima parte di questo libro, il Cansiglio è anche una delle aree che meriterebbero un'attenzione di ricerca più sistematica per l'arte rupestre, avendo già mostrato di conservare un archivio ambientale di qualità per il periodo che qui interessa.

\section{Il clima si stabilizza}

Dopo la fase fredda e le sue code glaciali, il clima entra progressivamente nella fase che dura, con oscillazioni minori, fino a oggi: l'Olocene, iniziato per convenzione a circa novemilasettecento anni prima dell'anno zero. È in questa fase lunga che si sviluppano le forme di vita umana che occuperanno il resto di questo libro: gruppi di cacciatori-raccoglitori sempre più adattati alla montagna, e poi, verso la fine del periodo che qui interessa, le prime comunità che iniziano a praticare l'agricoltura e l'allevamento.

Anche in questa fase lunga il clima non fu perfettamente costante. Il limite del bosco, nel corso dei millenni, ha conosciuto ripetuti abbassamenti temporanei, nell'ordine di sessanta, cento metri di quota, legati a oscillazioni climatiche minori. Più tardi, quando le comunità umane cominciano a praticare l'allevamento e il disboscamento sistematico, l'effetto delle attività umane sul limite del bosco diventa più marcato dell'effetto del clima stesso: un dettaglio che vale la pena tenere presente, perché segna il momento in cui il paesaggio alpino comincia a essere plasmato tanto dalla natura quanto dall'azione dell'uomo.

\begin{table}[htbp]
\centering
\small
\begin{tabular}{@{}p{3.4cm}p{2.6cm}p{5.4cm}@{}}
\toprule
\textbf{Fase} & \textbf{Anni prima dell'anno zero} & \textbf{Caratteristiche principali} \\
\midrule
Ultimo Massimo Glaciale & 19.000 & massima estensione dei ghiacciai alpini \\
Tardoglaciale, inizio & 17.000 & ritiro dei ghiacciai dagli anfiteatri pedemontani \\
Oscillazione di Ragogna / Gschnitz & 15.000--13.600 & battuta d'arresto della foresta, nuova avanzata glaciale \\
Bølling-Allerød & 12.700--11.000 ca. & foreste fino a 1.700--1.800 m, prime colonizzazioni stabili di quota \\
Laacher See Tephra & 11.050--10.950 & livello guida vulcanico, anticipa il Dryas Recente \\
Dryas Recente & 10.950--9.700 & clima secco e continentale, impatto sugli insediamenti non chiaro \\
Stadi Egesen/Kartell/Kromer & seconda metà del Dryas Recente, fino a inizio Olocene & nuova avanzata glaciale, datata localmente a 9.850 $\pm$ 1.100 (Val Cantone di Dosdé) \\
Olocene, inizio & 9.700 & stabilizzazione climatica, sviluppo delle culture mesolitiche \\
\bottomrule
\end{tabular}
\caption{Sequenza cronologica del Tardoglaciale e primo Olocene nelle Alpi italiane, secondo Ravazzi et al. (2007) salvo dove diversamente indicato. Le date sono espresse in anni prima dell'anno zero astronomico, ottenute sottraendo circa 1.950 anni dai valori cal BP originali, arrotondate al centinaio più vicino.}
\label{tab:cronologia}
\end{table}

\section{Cosa serve, di tutto questo}

Non tutti i dettagli di questa lunga vicenda climatica sono ugualmente utili per il problema pratico che questo libro vuole affrontare: capire dove cercare tracce di frequentazione preistorica sulle Alpi venete. Alcuni punti però sono decisivi, e vale la pena fissarli con chiarezza prima di proseguire.

Primo: prima di circa dodicimilasettecento anni prima dell'anno zero, cercare tracce di una presenza umana stabile in quota sulle Alpi venete è, con ogni probabilità, uno sforzo mal indirizzato. Il territorio non era ancora pronto ad accoglierla in modo continuativo, essendo ancora soggetto a un'intensa instabilità paraglaciale.

Secondo: da quel momento in poi, e per tutto il periodo che interessa questo libro, il limite superiore utile per la ricerca può essere fissato, con un margine di approssimazione dichiarato e accettabile, attorno ai milleseicento, millesettecento metri di quota nelle fasi più antiche, salendo eventualmente fino a millesettecentocinquanta, milleottocento nelle fasi più calde: sopra questa soglia, specialmente nelle fasi più antiche, la probabilità di trovare tracce di una frequentazione stabile cala sensibilmente.

Terzo: l'effetto del Dryas Recente sugli insediamenti resta un punto aperto, dichiarato tale dagli stessi specialisti del settore. Un modello di ricerca che assumesse un abbandono totale della montagna durante questa fase fredda, o al contrario una continuità totale, andrebbe oltre quanto i dati permettono di affermare.

Quarto, e forse più importante: non serve, per gli scopi di questo libro, una ricostruzione climatica di altissima precisione, calibrata anno per anno o quota per quota in ogni singola vallata. Le soglie riassunte in questo capitolo, onestamente dichiarate come approssimazioni, sono sufficienti per orientare la ricerca senza costruire un'illusione di precisione che i dati disponibili non permettono di sostenere.

Con questo quadro temporale fissato, il prossimo capitolo può affrontare la domanda successiva: chi viveva davvero su queste montagne, una volta diventate abitabili, e come si comportava.
